% ============================================================================
%  Semantic Runtime Specifications
%  An abstract machine and execution architecture for a declarative,
%  evidence-bearing, completion-directed language.
%
%  Self-contained: the object language is recapitulated from first principles;
%  only the standard external literature is cited.
% ============================================================================
\documentclass[11pt,a4paper]{article}

\usepackage[utf8]{inputenc}
\usepackage[T1]{fontenc}
\usepackage{lmodern}
\usepackage[a4paper,margin=2.4cm]{geometry}
\usepackage{amsmath,amssymb,amsthm,mathtools}
\usepackage{stmaryrd}
\usepackage{enumitem}
\usepackage{booktabs}
\usepackage{array}
\usepackage{xcolor}
\usepackage{microtype}
\usepackage[numbers,sort&compress]{natbib}
\usepackage[colorlinks=true,linkcolor=blue!45!black,citecolor=blue!45!black,
            urlcolor=blue!55!black]{hyperref}
\usepackage{cleveref}

% ---- theorem environments ----
\theoremstyle{plain}
\newtheorem{theorem}{Theorem}[section]
\newtheorem{lemma}[theorem]{Lemma}
\newtheorem{proposition}[theorem]{Proposition}
\newtheorem{corollary}[theorem]{Corollary}
\theoremstyle{definition}
\newtheorem{definition}[theorem]{Definition}
\newtheorem{invariant}[theorem]{Invariant}
\newtheorem{example}[theorem]{Example}
\theoremstyle{remark}
\newtheorem{remark}[theorem]{Remark}

% ---- operators and notation ----
\DeclareMathOperator{\dom}{dom}
\DeclareMathOperator{\cod}{cod}
\DeclareMathOperator{\emb}{embed}
\DeclareMathOperator*{\argmin}{arg\,min}
\DeclareMathOperator*{\argmax}{arg\,max}
\newcommand{\R}{\mathbb{R}}
\newcommand{\N}{\mathbb{N}}
\newcommand{\Unit}{[0,1]}
\newcommand{\Val}{\mathbb{V}}
\newcommand{\Conf}{\mathcal{C}}
\newcommand{\nor}{\oplus}
\newcommand{\Type}{\tau}
\newcommand{\env}{\rho}
\newcommand{\store}{\mu}
\newcommand{\cont}{\kappa}
\newcommand{\dbt}{\Delta}
\newcommand{\sys}{\Psi}
\newcommand{\mstate}{\varsigma}            % machine state
\newcommand{\stepm}{\longrightarrow}       % machine step
\newcommand{\evalto}{\Downarrow}           % big-step
\newcommand{\sem}[1]{\llbracket #1 \rrbracket}
\newcommand{\dist}{\mathsf{d}}
\newcommand{\target}{\Gamma_{\!f}}
\newcommand{\jspace}{\qquad}
\newcommand{\catp}{\varkappa}              % catalytic power
\newcommand{\Res}{\mathsf{R}}              % resolver
\newcommand{\Prov}{\mathsf{P}}             % provider
\newcommand{\reg}{\mathsf{Reg}}            % provider registry
\newcommand{\infrule}[3]{%
  \ensuremath{\displaystyle\frac{\,#2\,}{\,#3\,}}\;\textsc{(#1)}}
\newcommand{\rn}[1]{\textsc{#1}}

\title{\textbf{Semantic Runtime Specifications}\\[3pt]
\large An Abstract Machine and Execution Architecture for a Declarative,\\
Evidence-Bearing, Completion-Directed Language}

\author{Kundai Farai Sachikonye\\[2pt]
\small Technical University of Munich\\
\small \texttt{kundai.sachikonye@wzw.tum.de}}

\date{\today}

\begin{document}
\maketitle

\begin{abstract}
\noindent
We specify the runtime of a declarative language whose programs do not script a
sequence of steps but \emph{declare a completion condition} and accumulate
graded evidence toward it. Such a language poses runtime requirements that a
conventional evaluator does not meet: a sub-computation may be resolved by any
of several substrates (a deterministic primitive, an external tool, an oracle),
each carrying a confidence; the result a program seeks is specified
\emph{backward}, as a target rather than a procedure; and every value that
enters a conclusion must be \emph{grounded} in an observation rather than
freely generated. We give a complete, self-contained specification in five
parts. (i) A compact recapitulation of the object language---its values, a
confidence algebra on the unit interval, points and their resolution,
propositions and their argumentation semantics, and the iteration
constructs---so that the document stands alone. (ii) An \emph{abstract machine}
in the explicit-control--environment--store--continuation style, and a proof
that it is \emph{adequate}: it computes exactly the language's big-step
semantics. (iii) A \emph{resolver/provider architecture}: typed interfaces that
bind underspecified sub-computations to concrete substrates, with a
\emph{grounding} invariant that no value is committed to a conclusion without a
matching observation. (iv) A \emph{completion-directed} evaluation mechanism:
tasks are targets in a refinement hierarchy embedded in a metric space, reached
by backward navigation through \emph{virtual sub-states}; we prove the
logarithmic navigation bound and that the speed-up is impossible without virtual
sub-states. (v) A best-first \emph{scheduler} with a cascade of resolvers whose
end-to-end accuracy obeys a multiplicative law, together with a content-addressed
store and an observational decision log and state surface. The metatheory
establishes adequacy, machine safety, scheduler fairness and termination, the
cascade-accuracy bound and its saturation criterion, preservation of the
grounding invariant, and conservation of confidence within the unit interval.
The specification is precise enough to serve as an implementation blueprint and
is related throughout to the standard literature on abstract machines,
operational semantics, search, information geometry, and probabilistic
computation.

\medskip
\noindent\textbf{Keywords:} abstract machine, operational semantics, adequacy,
completion-directed evaluation, backward navigation, virtual sub-states,
resolver, scheduler, cascade, content-addressed store, grounding.
\end{abstract}

\tableofcontents
\bigskip\hrule

% ===========================================================================
\section{Introduction}
\label{sec:intro}
% ===========================================================================

\subsection{What an evidence-bearing, completion-directed runtime must provide}

A runtime is a precise account of how a program is executed: a machine that
turns syntax into results, together with the services---memory, scheduling,
dispatch---that the machine needs. For a conventional language the account is
well understood: an abstract machine in the tradition of \citet{landin1964,
felleisen1992} realises the language's operational semantics
\citep{plotkin1981,kahn1987,winskel1993}, and adequacy theorems
\citep{ager2003,reynolds1972} certify that the machine computes what the
semantics prescribes. The language whose runtime we specify here is not
conventional in three respects, each of which imposes a requirement that the
classical account does not address.

\begin{enumerate}[leftmargin=1.5em]
\item \textbf{Plural substrates with confidence.} A sub-computation---resolving
an ambiguous datum, scoring a hypothesis, answering a sub-query---may be
discharged by any of several \emph{substrates}: a deterministic primitive, an
external tool, or an oracle whose answer is uncertain. Each returns not just a
value but a \emph{confidence}. The runtime must dispatch uniformly to these
substrates and carry the confidence through.

\item \textbf{Backward specification.} A program declares \emph{what completion
looks like}---a target condition---rather than the steps to reach it. The
runtime must compute the trajectory to the target, and do so efficiently. We
show this is possible in logarithmic rather than linear work, but only through a
device, the \emph{virtual sub-state}, that the runtime must support.

\item \textbf{Grounding.} A conclusion must rest on observation. A value
produced by an oracle and never checked against the world is \emph{ungrounded};
the runtime must forbid such values from entering a committed result. This is a
liveness/safety requirement on the scheduler, not on the semantics.
\end{enumerate}

\subsection{Contributions}

We specify a runtime meeting these requirements and prove its guarantees.

\begin{itemize}[leftmargin=1.5em]
\item A self-contained recapitulation of the object language
(\cref{sec:language}), so the document needs no companion.
\item An abstract machine (\cref{sec:machine}) and an \emph{adequacy} theorem:
the machine's reflexive-transitive step relation computes exactly the big-step
semantics (\cref{thm:adequacy}).
\item A resolver/provider architecture (\cref{sec:resolvers}) with typed
interfaces and a \emph{grounding} invariant (\cref{inv:grounding}).
\item Completion-directed evaluation by backward navigation through virtual
sub-states (\cref{sec:completion}), with the logarithmic-navigation theorem
(\cref{thm:lognav}) and the \emph{necessity} of virtual sub-states
(\cref{thm:necessity}).
\item A best-first scheduler with a resolver cascade (\cref{sec:scheduler}),
the multiplicative accuracy law (\cref{thm:cascade}), and a content-addressed
store and observational project state (\cref{sec:store}).
\item The metatheory (\cref{sec:meta}): adequacy, machine safety, fairness and
termination, cascade accuracy and saturation, grounding preservation, and
confidence conservation.
\end{itemize}

\Cref{sec:trace} gives a worked execution; \cref{sec:related} situates the
design; \cref{sec:conclusion} concludes.

% ===========================================================================
\section{The Object Language, Recapitulated}
\label{sec:language}
% ===========================================================================

We fix a core calculus---the fragment the runtime must execute---together with
its semantic domains and big-step semantics. The presentation is compact but
complete: nothing below depends on external material beyond the cited standard
results.

\subsection{Core syntax}

\begin{definition}[Core terms]
\label{def:core}
Expressions $e$ and statements $s$ are
\begin{align*}
e &::= n \mid str \mid b \mid x \mid e\mathbin{op}e \mid e(\overline e)
       \mid e[e] \mid \mathsf{point}\{\overline{\ell:e}\}
       \mid \mathsf{resolve}(e,\varsigma)\\
s &::= \mathsf{item}\ x=e \mid x:=e \mid \mathsf{seq}(\overline s)
       \mid \mathsf{given}(e,s,s) \mid \mathsf{foreach}(x,e,s)
       \mid \mathsf{return}(e) \mid \mathsf{ensure}(e)\\
  &\quad\mid \mathsf{fun}(f,\overline x,s)
       \mid \mathsf{prop}(P,\overline M,s) \mid \mathsf{support}(M,e)
       \mid \mathsf{contradict}(M,e)\\
  &\quad\mid \mathsf{cycle}(x,e,s) \mid \mathsf{drift}(x,e,s)
       \mid \mathsf{roll}(s),
\end{align*}
with $op$ ranging over the arithmetic, comparison, and logical operators,
$\varsigma$ over resolution strategies (\cref{def:strategy}), and
$x,f,P,M,\ell$ over identifiers and labels. This is the fragment exercised by
the runtime; the remaining surface constructs reduce to it.
\end{definition}

\subsection{Semantic domains}

\begin{definition}[Values]
\label{def:values}
The value set $\Val$ is generated by
$v ::= \mathsf{num}(r)\mid\mathsf{str}(s)\mid\mathsf{bool}(b)\mid\mathsf{none}
\mid\mathsf{list}(\overline v)\mid\mathsf{clo}(\overline x,s,\env)
\mid\mathsf{pt}(p)\mid\mathsf{ent}(p)\mid\mathsf{floor}(F)
\mid\mathsf{corpus}(s)$,
where a \emph{point} $p=(\mathit{content},\mathit{conf}\in\Unit,
\overline{(m_i,p_i)})$ carries content, a confidence, and a sub-stochastic
list of interpretations ($\sum_i p_i\le1$); an \emph{entity} is a resolved
point; a \emph{floor} $F$ is a finite map from keys to weighted points; a
\emph{corpus} wraps a body of text.
\end{definition}

\begin{definition}[Confidence algebra]
\label{def:confalg}
Confidences combine by the \emph{noisy-or} on $\Unit$,
\[
c\nor d\;=\;1-(1-c)(1-d),
\]
which is associative, commutative, monotone, has unit $0$, and is bounded above
by $1$ (with $1$ absorbing). It is the probability that at least one of two
independent sources of confidence $c,d$ succeeds.
\end{definition}

\begin{definition}[Resolution strategy]
\label{def:strategy}
A \emph{resolution strategy} $\varsigma$ maps a point $p$ with normalised
interpretations $\overline{(m_i,\hat p_i)}$ to a determinate outcome: the
maximum-likelihood strategy returns the dominant interpretation when it clears
an acceptance threshold $\theta$ and otherwise an uncertain outcome carrying the
distribution and residual ambiguity $1-A(p)$, where $A(p)=H(p)/\log_2 k$ is the
normalised interpretation entropy; other strategies (Bayesian, conservative,
exploratory, full) re-weight or re-select accordingly. Resolving collapses
$\mathsf{pt}(p)$ to $\mathsf{ent}(\hat p)$.
\end{definition}

\begin{definition}[Environment, store, debate, system state]
\label{def:state}
An \emph{environment} $\env\colon\mathit{Var}\rightharpoonup\mathit{Loc}$ maps
variables to locations; a \emph{store} $\store\colon\mathit{Loc}\rightharpoonup
\Val$ maps locations to values (mutation updates $\store$). A \emph{debate
state} $\dbt\colon\mathit{Motion}\rightharpoonup(\mathcal{A}\times\mathcal{C})$
maps each motion to multisets of weighted affirmations and contentions. A
\emph{system state} $\sys=(\mathcal{H},\mathcal{F})$ pairs an append-only
\emph{decision log} $\mathcal{H}$ with a \emph{state surface} $\mathcal{F}$.
\end{definition}

\subsection{Big-step semantics (the specification to be realised)}

The semantics has expression judgements $\langle e,\env,\store\rangle\evalto
\langle v,\store'\rangle$ and statement judgements $\langle s,\env,\store,\dbt
\rangle\evalto\langle\chi,\store',\dbt'\rangle$ with completion $\chi\in\{\mathsf
{nrm}\}\cup\{\mathsf{ret}\,v\}$. We list the rules the runtime must honour; they
are standard call-by-value \citep{winskel1993,pierce2002} except for the three
distinctive forms.

\begin{gather*}
\infrule{B-Resolve}{\langle e,\env,\store\rangle\evalto\langle\mathsf{pt}(p),
  \store'\rangle\quad \hat p=\varsigma(p)}
  {\langle\mathsf{resolve}(e,\varsigma),\env,\store\rangle\evalto\langle
  \mathsf{ent}(\hat p),\store'\rangle}\\[8pt]
\infrule{B-Prop}{\dbt_0=[\overline{M\mapsto(\varnothing,\varnothing)}]\quad
  \langle s,\env,\store,\dbt_0\rangle\evalto\langle\_,\store',\dbt_1\rangle}
  {\langle\mathsf{prop}(P,\overline M,s),\env,\store,\dbt\rangle\evalto\langle
  \mathsf{nrm},\store'[P\mapsto\mathrm{verdict}_\varsigma(\dbt_1)],\dbt\rangle}\\
  [8pt]
\infrule{B-Support}{\langle e,\env,\store\rangle\evalto\langle\mathsf{num}(c),
  \store'\rangle\quad c'=\mathrm{clip}_\Unit(c)}
  {\langle\mathsf{support}(M,e),\env,\store,\dbt\rangle\evalto\langle\mathsf{nrm}
  ,\store',\dbt[M\,{+}{=}_{\mathcal A}(c',c',1)]\rangle}
\end{gather*}
where $\mathrm{verdict}_\varsigma$ scores each motion (an affirmations-minus-
contentions aggregate under strategy $\varsigma$, clipped to $\Unit$) and labels
it Supported/Contradicted/Inconclusive by thresholds. The iteration forms
$\mathsf{cycle}/\mathsf{drift}/\mathsf{roll}$ run their body once per floor
entry, per sampled corpus fragment, and until a convergence predicate fires,
respectively. We take this semantics as given and specify a machine that
realises it.

% ===========================================================================
\section{The Abstract Machine}
\label{sec:machine}
% ===========================================================================

We give a small-step abstract machine in the control--environment--store--
continuation style \citep{felleisen1992,ager2003}. The machine makes evaluation
order, the store, and the pending context explicit, and exposes exactly the
points at which the runtime dispatches to resolvers and providers
(\cref{sec:resolvers}).

\subsection{Configurations}

\begin{definition}[Machine state]
\label{def:mstate}
A \emph{machine state} is a tuple
\[
\mstate\;=\;\langle\,c,\ \env,\ \store,\ \cont,\ \dbt,\ \sys\,\rangle,
\]
where $c$ is the \emph{control} (a term being evaluated, or a value being
returned to a continuation), $\env$ the environment, $\store$ the store,
$\cont$ the \emph{continuation} (the pending context), $\dbt$ the debate state,
and $\sys$ the system state. Continuations are
\[
\cont ::= \mathsf{halt}\mid \mathsf{arg}(\overline e,\env,\cont)
\mid \mathsf{op}(op,v,\cont)\mid \mathsf{seq}(s,\env,\cont)
\mid \mathsf{bind}(x,\cont)\mid \mathsf{res}(\varsigma,\cont)\mid\cdots
\]
each frame recording one suspended elimination form. A state is \emph{final}
when $c$ is a value $v$ and $\cont=\mathsf{halt}$.
\end{definition}

\subsection{Transitions}

The step relation $\mstate\stepm\mstate'$ decomposes evaluation into atomic
steps. We give the representative rules; the rest follow the same discipline.

\begin{align*}
\langle e_1\mathbin{op}e_2,\env,\store,\cont,\dbt,\sys\rangle
  &\stepm \langle e_1,\env,\store,\mathsf{op}_1(op,e_2,\env,\cont),\dbt,\sys
  \rangle\\
\langle v_1,\env',\store,\mathsf{op}_1(op,e_2,\env,\cont),\dbt,\sys\rangle
  &\stepm \langle e_2,\env,\store,\mathsf{op}_2(op,v_1,\cont),\dbt,\sys\rangle\\
\langle v_2,\env,\store,\mathsf{op}_2(op,v_1,\cont),\dbt,\sys\rangle
  &\stepm \langle\sem{op}(v_1,v_2),\env,\store,\cont,\dbt,\sys\rangle\\[4pt]
\langle x,\env,\store,\cont,\dbt,\sys\rangle
  &\stepm \langle\store(\env(x)),\env,\store,\cont,\dbt,\sys\rangle\\[4pt]
\langle e(\overline e),\env,\store,\cont,\dbt,\sys\rangle
  &\stepm \langle e,\env,\store,\mathsf{arg}(\overline e,\env,\cont),\dbt,\sys
  \rangle\\
\langle\mathsf{clo}(\overline x,s,\env_0),\ldots,\mathsf{arg}(\overline v,
  \_,\cont),\ldots\rangle
  &\stepm \langle s,\ \env_0[\overline{x\mapsto\ell}],\ \store[\overline{\ell
  \mapsto v}],\ \cont,\ \dbt,\sys\rangle
\end{align*}
For the distinctive forms, the machine \emph{suspends to a resolver}: evaluating
$\mathsf{resolve}(e,\varsigma)$ pushes $\mathsf{res}(\varsigma,\cont)$, and when
the operand reduces to a point the machine emits a \emph{resolution request}
(\cref{sec:resolvers}) whose answer re-enters as the entity value; evaluating
$\mathsf{support}(M,e)$ updates $\dbt$ once $e$ is a number; evaluating
$\mathsf{prop}$ installs a fresh inner $\dbt_0$, runs the body, and writes the
verdict to the store and an entry to $\sys.\mathcal H$.

\begin{definition}[Evaluation by the machine]
\label{def:macheval}
Let $\stepm^{*}$ be the reflexive-transitive closure of $\stepm$. The machine
\emph{evaluates} a closed term $t$ to $v$ if
$\langle t,\varnothing,\store_0,\mathsf{halt},\dbt_0,\sys_0\rangle\stepm^{*}
\langle v,\_,\store',\mathsf{halt},\dbt',\sys'\rangle$.
\end{definition}

\subsection{Adequacy}

\begin{theorem}[Adequacy]
\label{thm:adequacy}
For every closed expression $e$ and store $\store$: the machine evaluates $e$ to
$v$ with final store $\store'$ \emph{iff} $\langle e,\varnothing,\store\rangle
\evalto\langle v,\store'\rangle$ in the big-step semantics; and likewise for
statements, with the machine reaching a final state carrying completion $\chi$,
store $\store'$, and debate $\dbt'$ iff $\langle s,\varnothing,\store,\dbt
\rangle\evalto\langle\chi,\store',\dbt'\rangle$. The machine is therefore a
correct and complete realisation of the semantics.
\end{theorem}
\begin{proof}[Proof sketch]
This is the standard correspondence between a big-step semantics and a
CEK/CESK machine \citep{ager2003,reynolds1972,felleisen1992}. ($\Leftarrow$,
soundness of the machine) By induction on the big-step derivation, exhibiting
for each rule a finite block of machine steps that pushes the appropriate
continuation frames, evaluates sub-terms (using the IH), and pops the frame with
the combined value; the store and debate threads are passed unchanged where the
rule does not touch them and updated exactly where it does (e.g.\ \rn{B-Support}
updates $\dbt$, \rn{B-Prop} installs and discharges $\dbt_0$). ($\Rightarrow$,
completeness) By induction on the length of $\stepm^{*}$, using a
\emph{decomposition lemma}: every non-final state uniquely decomposes into a
redex in an evaluation context (the continuation), so each maximal run
corresponds to a unique big-step derivation; the resolution-request steps
correspond to the premises of \rn{B-Resolve} with the resolver supplying the
$\varsigma(p)$ side condition. Determinacy of $\stepm$ (each state has at most
one successor) makes the correspondence a bijection on terminating runs.
\end{proof}

\begin{corollary}[Machine safety]
\label{cor:machinesafety}
On the statically typed fragment (no dynamic type), a well-typed initial state
never reaches a stuck non-final configuration: every non-final state has a
successor, except a deliberate abort from a failed $\mathsf{ensure}$ (a defined
outcome in the sense of \citet{hoare1969}) or a state awaiting a resolver
answer.
\end{corollary}
\begin{proof}
By adequacy (\cref{thm:adequacy}) the machine's reachable states are exactly the
sub-derivations of a big-step evaluation, and progress for the big-step
semantics on the static fragment \citep{wright1994,pierce2002} transfers: each
non-value control has a redex with a defined transition, the only non-progressing
outcomes being the defined abort and the resolver wait.
\end{proof}

% ===========================================================================
\section{Resolvers and Providers}
\label{sec:resolvers}
% ===========================================================================

The machine reduces a program to a sequence of \emph{requests}: primitive
operations to be computed, and underspecified sub-computations to be resolved.
We give the two interfaces that service them and the invariant that disciplines
their use.

\begin{definition}[Provider]
\label{def:provider}
A \emph{provider} for an operation name $op$ with signature
$\overline{\tau}\to\tau'$ is a (partial) function
$\Prov_{op}\colon\Val^{*}\rightharpoonup\Val$ that computes $op$ deterministically
on values of the declared types. A \emph{registry} $\reg$ is a finite map from
operation names to providers. The machine's $\mathsf{op}$ step consults $\reg$:
$\sem{op}=\reg(op)$.
\end{definition}

\begin{definition}[Resolver]
\label{def:resolver}
A \emph{resolver} is a map
\[
\Res\colon \mathit{Request}\times\mathit{Context}\longrightarrow
\Val\times\Unit,
\]
that answers an underspecified request (resolve a point, score a motion, satisfy
a sub-target) with a value and a confidence, given a context (the ambient store,
relevant evidence, and the active strategy). A resolver is classified by its
substrate: \emph{deterministic} (a pure function, e.g.\ applying a strategy to a
point's interpretation distribution), \emph{tool} (an external computation that
reads the world), or \emph{oracle} (a substrate whose answer is uncertain). Each
resolver carries a \emph{catalytic power} $\catp(\Res)\in\Unit$
(\cref{sec:scheduler}).
\end{definition}

\begin{definition}[Requests and observations]
\label{def:request}
A \emph{request} is \emph{generative} if it proposes a value (an oracle answer,
a candidate completion) and \emph{observational} if it reads the world (a tool
call, a measurement, a database lookup). A value $v$ produced by a generative
request is \emph{grounded} once an observational request has confirmed it within
tolerance; otherwise it is \emph{ungrounded}.
\end{definition}

\begin{invariant}[Grounding]
\label{inv:grounding}
No ungrounded value is committed to a conclusion. Operationally: a generative
resolution that contributes to a proposition verdict, a returned result, or a
published state must be matched, before commitment, by an observational
resolution that confirms it; the scheduler (\cref{sec:scheduler}) holds the
generative value in a pending state until the matching observation completes,
and discards it on failure.
\end{invariant}

\begin{remark}[Closure: action requires return]
\cref{inv:grounding} is the runtime's closure discipline: every outbound
(generative) step must be closed by an inbound (observational) step before its
effect is allowed to stand. Ungrounded generation---an oracle answer that never
faces the world---is precisely what the invariant forbids, and it is the
structural account of, and guard against, unfounded conclusions. The invariant
constrains the \emph{scheduler}, not the semantics: a program's meaning is
unchanged, but the runtime refuses to realise meanings built on uncorroborated
values.
\end{remark}

% ===========================================================================
\section{Completion-Directed Evaluation}
\label{sec:completion}
% ===========================================================================

A task declares a \emph{completion condition}---a target $\target$---rather than
a procedure. The runtime computes a trajectory that reaches $\target$. We show
this is done in logarithmic work by \emph{backward} navigation through a
refinement hierarchy embedded in a metric space, and that the speed-up depends
essentially on \emph{virtual sub-states}.

\subsection{The refinement hierarchy and its embedding}

\begin{definition}[Ternary refinement hierarchy]
\label{def:hierarchy}
A \emph{ternary refinement hierarchy} of depth $k$ is a sequence of partitions
$P_0,\dots,P_k$ of a target space with $|P_0|=1$, each block of $P_j$ refining
into exactly three children of $P_{j+1}$; the leaves $P_k$ number $N=3^{k}$. A
\emph{target} is a leaf; a \emph{trajectory} is a root-to-leaf path.
\end{definition}

\begin{definition}[Metric embedding]
\label{def:embedding}
An \emph{embedding} $\emb$ maps each node to a point of the cube $\Unit^{3}$ by
the base-3 expansion of its address along three coordinates, and equips
$\Unit^{3}$ with a product Fisher--Rao metric $\dist$
\citep{amari2016,cencov1982}. Under $\emb$, the parent of a node is obtained by
deleting the last base-3 digit of its coordinate---an $O(1)$ operation---and the
ball of radius $r$ about a leaf contains $O(r^{3}3^{k})$ leaves.
\end{definition}

\subsection{Virtual sub-states and backward navigation}

\begin{definition}[Virtual sub-state]
\label{def:virtual}
A \emph{sub-state decomposition} of a coordinate $s\in\Unit$ is a triple
$(s_1,s_2,s_3)\in\R^{3}$ with mean $\tfrac13(s_1+s_2+s_3)=s$. It is
\emph{virtual} if some $s_i\notin\Unit$ (a locally infeasible component whose
aggregate is feasible), and \emph{physical} otherwise.
\end{definition}

\begin{definition}[Backward navigation]
\label{def:backnav}
Given a target leaf $\target$, backward navigation returns the trajectory
$\target=p_k,p_{k-1},\dots,p_0$ by iterating the $O(1)$ parent map of
\cref{def:embedding}, each step admitting a (possibly virtual) sub-state
decomposition of the coordinate.
\end{definition}

\begin{theorem}[Logarithmic navigation]
\label{thm:lognav}
Backward navigation from any target leaf to the root of a depth-$k$ ternary
refinement hierarchy visits exactly $k=\log_3 N$ nodes, each in $O(1)$ work
under the embedding; total cost $O(\log_3 N)$.
\end{theorem}
\begin{proof}
Each step deletes one base-3 digit (the $O(1)$ parent map, \cref{def:embedding}),
moving up one level; from depth $k$ to $0$ takes $k$ steps; $k=\log_3 N$.
\end{proof}

\begin{theorem}[Necessity of virtual sub-states]
\label{thm:necessity}
Backward navigation restricted to \emph{physical} sub-state decompositions has
worst-case cost $\Theta(N)$---no better than forward enumeration. The
logarithmic bound of \cref{thm:lognav} requires virtual sub-states.
\end{theorem}
\begin{proof}[Proof sketch]
The geodesic short-cut in the continuous embedding requires evaluating the
metric at a resolution finer than the cell width of the next level, which is
attainable only through decompositions ranging over $\R^{3}$ rather than the
unit cube. Confined to physical decompositions, the navigator must instead
distinguish a node's child among the candidates at each level by direct
comparison; summing the candidate counts $\sum_{j=0}^{k}3^{j}=\Theta(N)$. Thus
the freedom to pass through locally infeasible (virtual) intermediate
coordinates is exactly what buys the exponential reduction; a forward
construction, which must realise each intermediate as a physical block, cannot
use it.
\end{proof}

\begin{corollary}[Observational forward, analytical backward]
\label{cor:asymmetry}
Forward construction realises each step physically and is $\Theta(N)$; backward
completion computes intermediates analytically and is $O(\log_3 N)$. A task is
therefore best executed by declaring its target and navigating backward, which
is the runtime's default for goal-resolution requests.
\end{corollary}

\subsection{Grounding meets completion}

Backward navigation produces a \emph{candidate} trajectory; by
\cref{inv:grounding} the leaf value it proposes is generative and must be observed before
commitment. Completion-directed evaluation thus interleaves an analytical
backward pass (cheap, generative) with an observational forward check
(grounding), and commits only when both agree within tolerance.

% ===========================================================================
\section{Scheduling, Cascades, and Determinism}
\label{sec:scheduler}
% ===========================================================================

The machine emits requests; the scheduler decides their order, dispatches them
to resolvers/providers, composes resolvers into cascades, and enforces grounding.

\subsection{Best-first scheduling}

\begin{definition}[Work item and priority]
\label{def:workitem}
A \emph{work item} is a pending request together with its current state
(coordinate) $s$ and its target $\target$. Its \emph{priority} is
$\pi=-\dist(s,\target)$: items closer to their target are scheduled first. The
scheduler maintains a priority queue of enabled items and repeatedly dispatches
the highest-priority one---best-first search in the sense of \citet{dijkstra1959,
hart1968}.
\end{definition}

\begin{definition}[Fair policy]
\label{def:fair}
A scheduling policy is \emph{fair} if every continuously enabled work item is
eventually dispatched. We require the scheduler to be fair, e.g.\ by ageing
priorities so that no enabled item is starved \citep{dijkstra1968}.
\end{definition}

\begin{theorem}[Liveness under fairness]
\label{thm:liveness}
Under a fair policy, every enabled resolution is eventually performed; hence if
a computation's request DAG is finite and acyclic, the computation terminates,
and if every leaf request has a resolver, it terminates with all requests
discharged.
\end{theorem}
\begin{proof}
Fairness gives that no enabled item waits forever (\cref{def:fair}); on a finite
acyclic DAG the dependency order is a well-founded partial order, so by
induction over it every item becomes enabled and is then, by fairness,
eventually dispatched; the process drains the DAG in finite time. The qualifier
on leaf resolvers ensures no item blocks for want of a substrate.
\end{proof}

\subsection{Resolver cascades and the accuracy law}

A request is often discharged by a \emph{cascade}: a sequence of resolvers each
refining the previous answer. We characterise the end-to-end accuracy.

\begin{definition}[Catalytic power and cascade]
\label{def:cascade}
The \emph{catalytic power} $\catp(\Res)\in\Unit$ of a resolver is the fraction
by which it reduces the residual distance of an answer to the truth: applying
$\Res$ to an answer at residual $\delta$ leaves residual $(1-\catp)\delta$. A
\emph{cascade} $\Gamma=\Res_1\diamond\cdots\diamond\Res_n$ applies the resolvers
in sequence.
\end{definition}

\begin{theorem}[Multiplicative accuracy law]
\label{thm:cascade}
The catalytic power of a cascade is
\[
\catp(\Res_1\diamond\cdots\diamond\Res_n)\;=\;1-\prod_{i=1}^{n}\big(1-
\catp(\Res_i)\big),
\]
the noisy-or aggregate (\cref{def:confalg}) of the stage powers; equivalently the
residual multiplies, $\delta_n=\delta_0\prod_i(1-\catp_i)$.
\end{theorem}
\begin{proof}
By induction on $n$. One stage: residual $\delta_0\mapsto(1-\catp_1)\delta_0$.
If after $n-1$ stages the residual is $\delta_0\prod_{i<n}(1-\catp_i)$, the
$n$-th leaves $\delta_0\prod_{i\le n}(1-\catp_i)$; the cascade power is
$1-\delta_n/\delta_0=1-\prod_i(1-\catp_i)$, which is the $\nor$-fold of the
$\catp_i$.
\end{proof}

\begin{corollary}[Saturation]
\label{cor:saturation}
A cascade drives the residual to $0$ (power to $1$) iff $\sum_i\catp(\Res_i)=
\infty$. A cascade of fixed per-stage power $\catp>0$ saturates geometrically;
one with $\catp_i=2^{-i}$ does not.
\end{corollary}
\begin{proof}
$\log\prod_i(1-\catp_i)=\sum_i\log(1-\catp_i)$ diverges to $-\infty$ iff
$\sum_i\catp_i=\infty$ (since $\log(1-x)\sim-x$), i.e.\ residual $\to0$; this is
the Borel--Cantelli dichotomy.
\end{proof}

\subsection{Concurrency and determinism}

\begin{proposition}[Determinism relative to the oracle assignment]
\label{prop:determinism}
Fix an assignment $\mathcal{O}$ of answers to every oracle request. Then the
runtime is deterministic: the machine step relation is a partial function
(\cref{thm:adequacy}), provider answers are functions of their arguments
(\cref{def:provider}), and oracle answers are fixed by $\mathcal{O}$. Two
independent work items (neither's target nor inputs depending on the other's
output) may be dispatched concurrently without affecting the result, because
their effects on disjoint store locations and distinct log entries commute.
\end{proposition}
\begin{proof}
Determinism: each source of choice is resolved---transitions by determinacy of
$\stepm$, providers by functionality, oracles by $\mathcal{O}$. Commutation:
independent items touch disjoint store locations (no read-write or write-write
hazard) and append to the log, an operation that is conflict-free up to the
recorded order \citep{lamport1978}; hence their two interleavings yield the same
final store and the same multiset of log entries.
\end{proof}

% ===========================================================================
\section{The Store and the Project State}
\label{sec:store}
% ===========================================================================

\subsection{Content-addressed store}

\begin{definition}[Content-addressed store]
\label{def:castore}
The runtime's long-term store is \emph{content-addressed}: the location of a
value is computed from the value, $\ell=\emb(\mathit{content})$, rather than
assigned independently. Retrieval is by coordinate proximity---the nearest
stored coordinate under $\dist$---and storage and retrieval share one operation.
There is no separate key space: an address is never supplied, only derived.
\end{definition}

\begin{proposition}[Storage--retrieval coincidence]
\label{prop:storeretrieve}
Under content addressing, $\mathrm{retrieve}(\emb(x))=x$ for any stored $x$
whose nearest neighbour is itself; storing $x$ and later retrieving by its
content recover $x$ exactly, and approximate content retrieves the nearest
stored item.
\end{proposition}
\begin{proof}
$\emb$ places $x$ at $\emb(x)$; retrieval returns $\argmin_{y}\dist(\emb(x),
\emb(y))$, which is $x$ when $x$ is stored and otherwise the nearest stored
coordinate. The embedding's injectivity on stored contents makes exact recovery
exact.
\end{proof}

\subsection{Decision log and state surface are observational}

\begin{definition}[Project threading]
\label{def:threading}
Two effects accompany evaluation beyond the store: a proposition verdict or a
resolution \emph{appends} a record to the decision log $\sys.\mathcal H$, and a
publish \emph{updates} the state surface $\sys.\mathcal F$. Both are write-only
with respect to the core: no machine transition reads $\mathcal H$ or
$\mathcal F$.
\end{definition}

\begin{theorem}[Observational separation]
\label{thm:observational}
Erasing $\mathcal H$ and $\mathcal F$ from every machine state yields exactly the
runtime over the bare core (store and debate only): the projection that drops
$\sys$ commutes with $\stepm$. Hence the log and surface record the computation
but do not steer it; the result is independent of them.
\end{theorem}
\begin{proof}
The only transitions touching $\sys$ append to $\mathcal H$ or update
$\mathcal F$ (\cref{def:threading}); none reads $\sys$. Therefore the projection
$\langle c,\env,\store,\cont,\dbt,\sys\rangle\mapsto\langle c,\env,\store,\cont,
\dbt\rangle$ maps each transition to the corresponding core transition and is a
simulation in both directions; the diagram commutes, and the final control,
store, and debate are unchanged by erasure.
\end{proof}

% ===========================================================================
\section{Metatheory}
\label{sec:meta}
% ===========================================================================

We collect the runtime's guarantees; several restate results proved above, here
assembled as the runtime's contract.

\begin{theorem}[Runtime contract]
\label{thm:contract}
The runtime satisfies:
\begin{enumerate}[label=(\roman*),leftmargin=1.9em,nosep]
\item \textbf{Adequacy} (\cref{thm:adequacy}): it computes exactly the object
language's big-step semantics.
\item \textbf{Safety} (\cref{cor:machinesafety}): on the static fragment no run
gets stuck except at a defined abort or a resolver wait.
\item \textbf{Liveness} (\cref{thm:liveness}): under fair scheduling every
enabled resolution is performed; finite acyclic request DAGs terminate.
\item \textbf{Efficiency} (\cref{thm:lognav}): goal resolution by backward
navigation is $O(\log_3 N)$, optimal up to the necessity of virtual sub-states
(\cref{thm:necessity}).
\item \textbf{Accuracy} (\cref{thm:cascade}, \cref{cor:saturation}): a resolver
cascade has end-to-end power $1-\prod_i(1-\catp_i)$ and saturates iff
$\sum_i\catp_i=\infty$.
\item \textbf{Grounding} (\cref{inv:grounding}): no ungrounded value is
committed to a conclusion.
\item \textbf{Confidence conservation}: every confidence the runtime computes
lies in $\Unit$.
\item \textbf{Observational state} (\cref{thm:observational}): the decision log
and state surface record but do not alter the computation.
\end{enumerate}
\end{theorem}
\begin{proof}
(i)--(vi), (viii) are the cited theorems. (vii): confidences arise from resolver
answers (in $\Unit$ by \cref{def:resolver}), from clipping in \rn{B-Support},
from strategy outputs (in $\Unit$ by \cref{def:strategy}), and from cascade
aggregation by $\nor$, which is bounded above by $1$ and below by $0$
(\cref{def:confalg}); all constructors preserve $\Unit$, so by induction over
the evaluation every computed confidence is in $\Unit$.
\end{proof}

\begin{theorem}[Grounding preservation]
\label{thm:groundpres}
If the scheduler holds every generative value pending until a matching
observational resolution succeeds, then \cref{inv:grounding} holds of every
reachable state: no reachable committed conclusion depends on an ungrounded
value.
\end{theorem}
\begin{proof}
By induction on the run. Initially no value is committed. The only steps that
commit a value to a conclusion are: writing a proposition verdict (\rn{B-Prop}),
returning a result, and publishing state. By hypothesis the scheduler does not
enable these on a value still pending grounding; thus at each such step the
contributing value has a completed matching observation, preserving the
invariant. Generative values whose observation fails are discarded and never
reach a committing step.
\end{proof}

\begin{remark}[Termination is conditional]
As for any language with general recursion and unbounded iteration
\citep{hopcroft2006}, termination of $\mathsf{while}$/$\mathsf{roll}$ is
undecidable; \cref{thm:liveness} gives termination for finite acyclic request
DAGs and bounded iteration. The runtime therefore guarantees progress and
fairness unconditionally and termination under the stated finiteness condition,
exactly as the object language guarantees termination only for its bounded
constructs.
\end{remark}

% ===========================================================================
\section{A Worked Execution}
\label{sec:trace}
% ===========================================================================

Consider resolving an ambiguous datum inside a proposition, with a two-stage
resolver cascade and the grounding check.

\begin{enumerate}[leftmargin=1.5em]
\item The machine evaluates $\mathsf{prop}(P,\overline M,s)$ (\rn{B-Prop}):
installs $\dbt_0$, runs the body.
\item The body reaches $\mathsf{item}\ d=\mathsf{resolve}(\mathsf{point}\{
\dots\},\varsigma)$. The machine builds $\mathsf{pt}(p)$ and pushes
$\mathsf{res}(\varsigma,\cont)$, emitting a \emph{resolution request} for $p$.
\item The scheduler dispatches the request to a cascade $\Res_1\diamond\Res_2$:
$\Res_1$ a deterministic strategy resolver ($\catp_1=0.6$) producing a candidate
entity; $\Res_2$ an oracle ($\catp_2=0.5$) refining it. By \cref{thm:cascade}
the cascade power is $1-(0.4)(0.5)=0.8$.
\item The oracle answer is \emph{generative}; by \cref{inv:grounding} it is held
pending. The scheduler issues an \emph{observational} request (a tool lookup);
on agreement within tolerance the entity is grounded and committed as $d$, with
confidence the cascade aggregate; on disagreement it is discarded and the next
candidate sought.
\item Back in the proposition body, $\mathsf{given}\ \phi:\ \mathsf{support}(M,
d.\mathit{conf})$ updates $\dbt_0$ (\rn{B-Support}). After the body,
$\mathrm{verdict}_\varsigma(\dbt_1)$ is written to the store and appended to the
decision log $\sys.\mathcal H$ (\cref{def:threading}).
\end{enumerate}
By \cref{thm:adequacy} the value computed equals the big-step result; by
\cref{thm:observational} the log entry does not change it; by
\cref{thm:groundpres} the committed verdict rests only on grounded values.

% ===========================================================================
\section{Related Work}
\label{sec:related}
% ===========================================================================

\paragraph{Abstract machines and adequacy.} The machine is in the SECD/CEK/CESK
lineage \citep{landin1964,felleisen1992}; the adequacy proof follows the
functional correspondence between evaluators and machines
\citep{ager2003,reynolds1972} and the standard relationship to big-step and
structural operational semantics \citep{kahn1987,plotkin1981,winskel1993}. Type
safety transfers by the syntactic method \citep{wright1994,pierce2002}.

\paragraph{Dispatch and effects.} The resolver/provider split is a typed
dispatch interface in the spirit of effect handlers and the monadic account of
computational effects \citep{moggi1991,wadler1992}; the registry is a small
module system \citep{cardelli1985}.

\paragraph{Search and navigation.} Best-first scheduling by distance-to-target
is Dijkstra/A* search \citep{dijkstra1959,hart1968}; the backward-navigation
bound and the metric embedding draw on information geometry
\citep{amari2016,cencov1982} and a fixed-point view of convergent refinement
\citep{banach1922}. The content-addressed store realises storage and retrieval
as one coordinate operation.

\paragraph{Uncertainty.} Confidences and their noisy-or aggregation place the
runtime in the probabilistic-computation tradition \citep{kozen1981,mciver2005}
and the theory of evidence \citep{shafer1976}; proposition verdicts realise a
quantitative argumentation \citep{dung1995}.

\paragraph{Concurrency.} Fairness and the conflict-free composition of
independent work items follow classical concurrency \citep{dijkstra1968,
lamport1978}.

The novelty is the assembly: an adequate machine whose request stream is
serviced by substrate-agnostic resolvers under a grounding invariant, with goal
resolution by backward navigation through virtual sub-states, a cascade-accuracy
law, and an observational project state---a runtime tailored to a language that
declares targets and reasons under evidence.

% ===========================================================================
\section{Conclusion}
\label{sec:conclusion}
% ===========================================================================

We have specified the runtime of a declarative, evidence-bearing,
completion-directed language and proved its contract. An abstract machine
realises the language's semantics adequately and safely; a resolver/provider
architecture binds underspecified sub-computations to deterministic, tool, and
oracle substrates while a grounding invariant forbids any ungrounded value from
entering a conclusion; goal resolution proceeds by backward navigation through a
metric-embedded refinement hierarchy in logarithmic work, provably impossible
without virtual sub-states; a fair best-first scheduler composes resolvers into
cascades whose accuracy follows a multiplicative law with a sharp saturation
criterion; and a content-addressed store with an observational decision log and
state surface completes the picture. The specification is self-contained and
constructive: the machine fixes the evaluator, the interfaces fix the dispatch,
the navigation fixes goal resolution, and the scheduler fixes the order and the
grounding---each with a stated and proved guarantee, so that an implementation
has both a blueprint and a contract to meet.

% ===========================================================================
\bibliographystyle{plainnat}
\bibliography{references}
% ===========================================================================
\end{document}
